\section{Thảo luận}\label{sec:disc}
\textbf{Vai trò của đồ thị trong truy hồi.} So với mốc CLIP + adapter, QA-Graph tăng MRR trên toàn bộ ảnh test ở cả VG$\cap$COCO và PSG$\cap$COCO; xáo cạnh làm kết quả giảm xuống dưới mốc này. Tuy nhiên, phần tăng tập trung ở nhóm truy vấn trực tiếp. Cách gán nhóm gián tiếp dựa trên việc thiếu nhãn đối tượng trong ảnh chuẩn, nên nhóm này còn trộn lẫn chú thích thiếu, từ đồng nghĩa và nhu cầu suy luận quan hệ. Kết quả chưa tách được riêng đóng góp của suy luận nhiều bước. Việc cạnh quan hệ dày có thể làm nhãn khó phân biệt và việc CLIP đã biểu diễn tốt đối tượng trong ảnh là các giả thuyết phù hợp với ablation hiện có, chưa phải nguyên nhân được xác nhận độc lập.

\textbf{Giải thích quyết định của mô hình.} Các phép xóa cạnh cho thấy một số ít cạnh đối tượng có tác động rõ lên điểm ảnh đứng đầu. Fidelity của giải thích đối với hàm chấm điểm không đồng nghĩa với validity thị giác: cạnh có thể đúng trong đồ thị nhưng chú thích đối tượng vẫn thiếu hoặc sai; đường đi ngắn cũng không chứng minh quan hệ được mô hình hiểu theo nghĩa của câu truy vấn. Số đường đi quan hệ dài hai cạnh rất ít trong lượt khảo sát. Việc ảnh chuẩn tăng hạng sau khi xóa cạnh cũng không phải kết quả tất yếu, vì điểm của toàn bộ ảnh ứng viên được tính lại. Phân tích ảnh cho thấy một ảnh khác ID vẫn có thể phù hợp với caption, do đó cần trình bày riêng khớp ID, mức phù hợp ngữ nghĩa và tác động của cạnh lên điểm.

\textbf{Mốc so sánh.} Adapter học được tạo phần đáng kể trong mức tăng so với CLIP giữ nguyên trọng số. Baseline khớp tên đối tượng là đối chứng quan trọng cho phương pháp có quan hệ: AdaMF-MAT tái lập được trên MKG-W và MKG-Y, nhưng khi chuyển sang xếp hạng ảnh không vượt baseline đơn giản này. Kết quả ablation và đối chứng XAI cho thấy tín hiệu cấu trúc có ảnh hưởng, đồng thời giới hạn kết luận ở đúng phép đo và tập truy vấn đã khảo sát.
